\documentclass{article}

\title{Ch. 11 CW}
\author{Principles of Microeconomics}
\date{\today}

\begin{document}

\maketitle

\subsection*{Instructions} Please justify all of your answers and submit your solutions on Gradescope by the beginning of next class

\subsection*{Problems}

\begin{enumerate}

\item Four roommates are planning to spend the weekend in their dorm room watching old movies, and they are debating how many to watch. Here is their willingness to pay for each film:

\begin{center}
\begin{tabular}{| c | c | c | c | c |}
\hline
& Dwayne & Javier & Salman & Chris \\
\hline
First film & \$7 & \$5 & \$3 & \$2 \\
\hline 
Second film & 6 & 4 & 2 & 1 \\
\hline 
Third film & 5 & 3 & 1 & 0 \\
\hline 
Fourth film & 4 & 2 & 0 & 0 \\
\hline 
Fifth film & 3 & 1 & 0 & 0 \\
\hline
\end{tabular}
\end{center}

	\begin{enumerate}
	
	\item Within the dorm room, is the showing of a movie a public good? Why or why not?
	
	\item If it costs \$8 to stream a movie, how many movies should the roommates stream to maximize total surplus?
	
	\item If they choose the optimal number from part (b) and then split the cost of streaming the movies equally, how much surplus does each person obtain from watching the movies? 
	
	\item Is there any way to split the cost to ensure that everyone benefits? What practical problems does this solution raise?
	
	\item Suppose they agree in advance to choose the efficient number and to split the cost of the movies equally. When Dwayne is asked his willingness to pay, will he have an incentive to tell the truth? If so, why? If not, what will he be tempted to say?
	
	\item What does this example teach you about the optimal provision of public goods?
	
	\end{enumerate}
	
\item Some economists argue that private firms will not undertake the efficient amount of basic scientific research.

	\begin{enumerate}

	\item Explain why this might be so. In your answer, classify basic research as one of the four types of goods.
	
	\item What sort of policy has the United States adopted in response to this problem?
	
	\item It is often argued that this policy increases the technological capability of American producers relative to that of foreign firms. Is this argument consistent with your classification of basic research in part (a)? (Hint: Can excludability apply to some potential beneficiaries of a public good and not others?)
	
	\end{enumerate}
	
\item Two towns, each with three residents, are deciding whether to put on a fireworks display to celebrate the New Year. Fireworks cost \$360. In each town, some people enjoy fireworks more than the others.

	\begin{enumerate}
	
	\item In the town of Bayport, each of the residents values the public good as follows:
	
	\begin{center}
	\begin{tabular}{| c | c |}
	\hline
	Frank & \$50 \\
	\hline
	Joe & \$100 \\
	\hline
	Callie & \$300 \\
	\hline
	\end{tabular}
	\end{center}
	
	Would fireworks pass a cost-benefit analysis? Explain. 
	
	\item The mayor of Bayport proposes to decide by majority rule and, if the fireworks referendum passes, to split the cost equally among all residents. Who would vote in favor, and who would vote against? Would the vote yield the same answer as the cost-benefit analysis?
	
	\item In the town of River Heights, each of the residents values the public good as follows:
	
	\begin{center}
	\begin{tabular}{| c | c |}
	\hline
	Nancy & \$20 \\
	\hline
	Bess & \$140 \\
	\hline
	Ned & \$160 \\
	\hline
	\end{tabular}
	\end{center}
	
	Would fireworks pass a cost-benefit analysis? Explain.
	
	\item The mayor of River Heights also proposes to decide by majority rule and, if the fireworks referendum passes, to split the cost equally among all residents. Who would vote in favor, and who would vote against? Would the vote yield the same answer as the cost-benefit analysis?
	
	\item What do you think these examples say about the optimal provision of public goods?
		
	\end{enumerate}

\item Many transportation systems, such as the Washington, D.C., Metro (subway), charge higher fares during rush hours than during the rest of the day. Why might they do this?

\item Four roommates are thinking about hiring a housekeeper. The cost of providing an hour of housekeeping is constant at \$20. One roommates demand for housekeeping is given by $Q_D = 5 - \frac{1}{2}P$ where $Q_D$ is denoted in hours. The other three roommate's individual demand for housekeeping is given by $Q_D = 3 - \frac{1}{2}P$.

	\begin{enumerate}
	
	\item What is the optimal number of hours of housekeeping that the apartment should buy?
	
	\item How many hours of housekeeping would each of the roommates buy if they all lived alone? 
		
	\end{enumerate} 

\end{enumerate}

\end{document}